\chapter{Traces, credit and mechanistic evidence}
A trace says that a component ran. A gradient says that an output is locally
sensitive to a variable. A credit rule assigns a share of a defined value.
An intervention asks what changes when a mechanism is deliberately altered.
These are different questions, even when their answers are drawn on the same
graph. Treating one as another can make a convincing-looking explanation wrong.

Chapters 15 and 16 gave executions semantics and resource ledgers. This
chapter gives explanatory claims a similarly explicit contract. We begin
with a synthetic system whose equations are known, implement interventions,
and examine where simple attribution rules fail. No trained Evolutor model
or experimentally discovered biological circuit is claimed here.

\section{Observation can fit perfectly and still mislead}
Let one externally assigned scalar $u$ determine four internal quantities:
\begin{equation}
 a=2u,\qquad b=u,\qquad z=u,\qquad y=a+3b.
 \label{eq:scm}
\end{equation}
The variable $z$ is logged but never used to compute $y$. On every unmodified
run, however, $y=5z$. A predictor using only $z$ can therefore fit these
observations exactly. It has learned a predictive relationship, not a
mechanistic dependency of $y$ on $z$.

\Plate{01-causal}{The declared structural mechanism has a common cause
$u$, two paths into $y$, and a bystander $z$. An intervention replaces one
equation and cuts its incoming dependency. It does not redraw the remaining
equations merely to preserve correlations seen in the original trace.}{fig:causal}

A \Term{structural intervention} $\operatorname{do}(a=0)$ replaces $a=2u$
with $a=0$ while leaving $b=u$ and $y=a+3b$ intact. Descendants must then
be recomputed. At $u=1$, the unmodified result is 5 and the intervention
result is 3. Replacing the saved value of $a$ in a log after execution would
leave the saved $y$ at 5; that is an edited record, not the declared intervention.

\Listing{evaluate}

The evaluator follows topological order. A supplied intervention value wins
over the corresponding equation; other equations use the resulting parent
values. Unknown targets and nonfinite values are rejected. Its generated
outputs at $u=1$ are
\begin{center}\input{\Generated/interventions.tex}\end{center}
Intervening on $z$ changes $z$ but not $y$. Intervening directly on $y$
is also allowed by this model: it replaces the output equation itself.
That intervention answers a different question from modifying one of its
causes. The target and boundary of an intervention belong in the report.

\section{Necessity depends on what else remains available}
Let two binary components $A,B$ produce $y=A\lor B$. With both active,
ablating either one leaves the output at one. It would be wrong to conclude
that both components are irrelevant. They are redundant for this input and
output criterion. Ablating both changes the answer.

Now let $y=A\land B$. Either single ablation changes the output from one
to zero, but neither component alone is sufficient. These mechanisms have
the same unmodified output and different intervention tables.

\Plate{02-effects}{Complete binary response tables expose redundancy and
synergy. Rows vary $A$, columns vary $B$. A single ablation from the upper-right
jointly active case cannot characterize all four conditions. Logical response
tables are model interventions, not biological knock-out measurements.}{fig:effects}

For a declared value function $v(S)$, with only components in coalition $S$
enabled, define the two-component interaction
\begin{equation}
 I_{AB}=v(\{A,B\})-v(\{A\})-v(\{B\})+v(\varnothing).
\end{equation}
It is $-1$ for OR and $+1$ for AND. For $v(S)=2\,1_{A\in S}+3\,1_{B\in S}$,
it is zero. The sign describes departure from additivity under this
coalition definition and baseline, not an intrinsic property of a component
independent of task or context.

In a neural system, disabling a component also needs a definition: zeroing,
mean replacement, resampling from a reference distribution and substituting
another example need not have the same effect. Each may create states not
encountered during training. A large effect can reveal sensitivity to the
intervention artifact rather than the intended feature.

\section{Derive credit as an average over ordering}
Suppose we want to allocate the change $v(N)-v(\varnothing)$ among players
in finite set $N$. For an ordering $\pi$, a player's marginal contribution
is the value gained when it joins the players preceding it. Average over
all orderings:
\begin{equation}
 \phi_i=\frac{1}{|N|!}\sum_{\pi}
 \left[v(P_i^\pi\cup\{i\})-v(P_i^\pi)\right].
 \label{eq:shapley}
\end{equation}
This is a \Term{Shapley allocation} for the specified game. Along each
ordering, the marginal terms telescope from the empty to the full coalition.
Therefore $\sum_i\phi_i=v(N)-v(\varnothing)$. The efficiency property
follows algebraically; it does not establish that the chosen game is causal.

\Plate{03-shapley}{Two orderings allocate OR's unit value differently.
Averaging assigns one half to each player. For AND the second arriving player
receives the unit increment instead; averaging again gives one half each.
Identical credit can hide different mechanisms.}{fig:shapley}
\Listing{shapley}

The reference requires unique player names and a nonempty player set. It
assumes a deterministic, pure value function: the same coalition must mean
the same intervention and evaluation. It enumerates permutations, so runtime
is factorial rather than a scalable attribution service. Exact enumeration
is useful here because it provides an inspectable oracle for small examples.

\begin{center}\input{\Generated/credit.tex}\end{center}
OR and AND both give $(0.5,0.5)$, although their necessity and sufficiency
patterns differ. The additive mechanism gives $(2,3)$. A constant baseline
of seven changes none of these marginal credits because it cancels in every
difference. Reporting credits without reporting the coalition semantics
throws away the information needed to interpret them.

The tests independently compute the equivalent subset-weighted expression
\[
 \phi_i=\sum_{S\subseteq N\setminus\{i\}}
 \frac{|S|!(|N|-|S|-1)!}{|N|!}\,[v(S\cup\{i\})-v(S)].
\]
The weight counts how many orderings place exactly $S$ before $i$.
Agreement between these implementations, plus efficiency and symmetry
checks, guards against indexing mistakes. It does not validate a real
system's intervention distribution.

\section{A trace through time needs timed interventions}
Consider a scalar memory with $h_0=0$:
\begin{equation}
 h_t=\lambda h_{t-1}+x_t,\qquad 0\leq\lambda\leq1.
\end{equation}
Replace input $x_k$ by $x'_k$ while holding all other inputs fixed and
rerunning all later transitions. Let $\Delta_t=h_t-h'_t$. Before $k$ the
runs agree; at $k$, $\Delta_k=x_k-x'_k$; afterward,
$\Delta_t=\lambda\Delta_{t-1}$. Thus
\begin{equation}
 \Delta_t=(x_k-x'_k)\lambda^{t-k},\qquad t\geq k.
 \label{eq:decay}
\end{equation}
This is an exact finite intervention result for this linear mechanism, not
a general identity for nonlinear recurrent models.

\Plate{04-time}{A patch to the second input propagates forward through
the unrolled recurrence. The generated curves use four unit inputs and
$\lambda=0.5$. Earlier states remain unchanged; later differences decay
by one factor of $\lambda$ per transition.}{fig:time}
\Listing{memory}

The code uses zero-based patch indices: \texttt{patch=(1,0)} changes the
second input. The generated table uses one-based times:
\begin{center}\input{\Generated/memory.tex}\end{center}
Patching $h_k$ directly is not the same interface as patching $x_k$.
An input patch computes $h'_k=\lambda h_{k-1}+x'_k$; a state patch can
assign any admissible $h'_k$. At later times both differences follow the
same scalar propagation only because this model is linear and subsequent
inputs are fixed. If future inputs depend on the changed state, the closed
loop must be recomputed too.

With $\lambda=0$, effects vanish after one further step. With $\lambda=1$,
they persist unchanged. A gradient $\partial h_t/\partial x_k=\lambda^{t-k}$
matches the finite effect divided by its magnitude here. In a nonlinear
system, gradients are local and can miss threshold crossings, saturation
or changed routes. Agreement in this example is a consequence of its
equations, not a license to equate gradients with counterfactuals everywhere.

\section{Test explanations on interventions, not only observations}
Return to Equation~\ref{eq:scm}. Surrogate $\hat y=5z$ has zero error on
all unmodified runs. At $u=1$, $\operatorname{do}(z=0)$ makes its prediction
zero while the true output remains five. Intervening on $a$ or $b$ also
breaks the surrogate's prediction. A perfect observational fit can therefore
fail the simplest mechanistic tests.

\Plate{05-evidence}{An explanation should be evaluated under a declared
distribution of interventions. The bystander surrogate fits the observational
line exactly but fails when the input mechanism is surgically changed.
Intervention coverage is additional evidence, not implied by trace agreement.}{fig:evidence}

A practical evaluation separates development interventions from held-out ones.
Declare the output metric and tolerance, intervention targets and magnitudes,
input distribution, random seeds, and invalid-state policy. Report failures
by intervention family rather than only one average. A surrogate can reproduce
single-node patches and still fail joint patches or long-horizon effects.
Keeping a list of all failures makes the explanation falsifiable.

Current circuit-tracing work provides a useful comparison: replacement graphs
can support interventions, but validation must examine whether those
interventions match the underlying model. The 2025 methods report discusses
fixed attention patterns, reconstruction errors and limits of replacement-model
faithfulness \cite{circuits}. Our synthetic evaluator avoids those approximation
issues by declaring exact equations; that makes it a testing lesson, not an
implementation or validation of that research system.

\section{An evidence contract for Evolutor}
For each proposed explanation, store the executable mechanism or exact model
revision, the selected events, the intervention operator, the recomputation
boundary, the response metric and the test set. A trace hash can bind what
was observed; it cannot prove the trace is complete or that an unobserved
component had no effect. A successful replay verifies reproducibility under
the replay contract, not the truth of every causal interpretation attached to it.

Keep explanatory claims at their earned level. A route log supports which
expert was selected. A zero-ablation measures sensitivity to removing its
output in that context. A controlled replacement can test a more specific
functional hypothesis. None alone establishes biological equivalence,
general task competence or a complete decomposition of learned computation.

\section{Exercises with worked answers}
\begin{enumerate}
\item Compute $\operatorname{do}(a=4)$ at $u=1$. \textbf{Answer:}
$b=1$, so $y=4+3=7$. The equation for $a$ is replaced, not added to.
\item Set both $a$ and $b$ to zero. \textbf{Answer:} $y=0$ regardless
of $z$. Descendant recomputation uses both replacements.
\item Why is $z$ predictive? \textbf{Answer:} It shares cause $u$ with
the useful paths. On observations $y=5u=5z$, without a $z\to y$ edge.
\item Does zeroing a log field implement a causal patch? \textbf{Answer:}
No; dependent computations must rerun under the replaced equation.
\item Derive OR's interaction. \textbf{Answer:} $1-1-1+0=-1$.
Either component substitutes for the other in this game.
\item Derive AND's interaction. \textbf{Answer:} $1-0-0+0=1$.
The joint contribution exceeds the sum of singleton contributions.
\item Explain identical OR and AND Shapley credits. \textbf{Answer:}
Each is symmetric with total value one. Averaging over orderings splits
that total equally despite different response tables.
\item Add seven to every coalition value. \textbf{Answer:} Each marginal
difference is unchanged; credits allocate the change above the baseline.
\item Patch the second unit input to zero at $\lambda=0.5$.
\textbf{Answer:} State differences at times one through four are
$0,1,0.5,0.25$, matching Equation~\ref{eq:decay}.
\item Can the patch change time one? \textbf{Answer:} Not in this causal
recurrence. Such a result would indicate indexing or replay leakage.
\item Construct a failed surrogate test. \textbf{Answer:} At $u=1$,
set $z=0$. True $y$ remains five; $5z$ predicts zero despite perfect
observational fit.
\item Design a stronger evaluation than single-node ablation.
\textbf{Answer:} Hold out inputs, paired interventions, patch magnitudes
and temporal horizons; compare predicted and actual response changes with
declared tolerances and report where the explanation fails.
\end{enumerate}

\section{Vocabulary and the next question}
\Term{Counterfactual execution}: a rerun under a specified altered mechanism
with the other declared conditions fixed. \Term{Redundancy}: alternative
components can preserve a chosen output. \Term{Synergy}: a joint effect
exceeds the additive comparison under a specified baseline.
\Term{Mechanistic faithfulness}: agreement of an explanation with the system's
responses over the tested interventions, not merely agreement on predictions.

The next chapter defines a DNA-addressed state primitive for DOGMA research.
It will earn simple claims through invariants, traces and counterexamples
before asking whether a larger trained architecture can make useful use of it.
